% Part: normal-modal-logic
% Chapter: tableaux
% Section: countermodels
%READERNOTE{SOL6-A181: निर्णय प्रक्रियेतील नवे पूर्वचिन्ह देणारा नियम प्रत्येक पूर्वचिन्हयुक्त चिन्हांकित सूत्रासाठी एकदा लावायचा असतो; वापरलेली पूर्वचिन्हे घेणाऱ्या नियमाची व्याप्ती त्याच्या तात्काळ विस्तारांपुरती आहे. थांबण्याचा सांत इनपुटवरील पुरावा स्पष्ट केला. मनमानी गृहीतकसंचासाठी संपूर्णतेचा निष्कर्ष जपून, प्रत्यक्ष सांत निर्णय प्रक्रियेचे विधान सांत गृहीतकसंचापुरते केले. उदाहरणांतील वापरलेली पूर्वचिन्हे 1.n या आकाराची आहेत हे स्पष्ट केले; सर्व सहा टॅब्लो आणि दोन्ही प्रतिमाने जपली.}

\documentclass[../../../include/open-logic-section]{subfiles}

\begin{document}

\olfileid{nml}{tab}{cou}

\olsection{\usetoken{P}{tableau} मधून प्रतिदृष्टांत प्रतिमाने}

\begin{explain}
  संपूर्णता प्रमेयाचा पुरावा
  $\Entails !A$ असेल तर $\Proves !A$
  हे दाखवतोच; शिवाय $\Entails/ !A$
  असेल, तर~$!A$ साठी प्रतिदृष्टांत
  प्रतिमान बांधण्याची पद्धतही देतो.
  $\Log{K}$ साठी ही पद्धत
  \emph{निर्णय प्रक्रिया} आहे. कारण
  $\Entails/ !A$ असे समजू. मग
  \olref[cpl]{prop:complete-tableau} चा
  पुरावा संपूर्ण !!{tableau} बांधण्याची
  पद्धत देतो. ती नेहमी थांबतेही.
  $\Log{K}$ चे विधानीय नियम कमी गुंतागुंतीची
  पूर्वचिन्हयुक्त !!{formula}s जोडतात;
  म्हणजे प्रत्येक
  $\sFmla{S}{!A}[\sigma]$ या चिन्हांकित
  सूत्रासाठी प्रत्येक विधानीय नियम
  शाखेवर एकदाच लावावा लागतो.
  नवी पूर्वचिन्हे फक्त
  \iftag{prvBox}{$\TRule{\False}{\Box}$}{}\iftag{notprvBox,notprvDiamond}{}{
    आणि }\iftag{prvDiamond}{$\TRule{\True}{\Diamond}$}{}
  \iftag{notprvBox,notprvDiamond}{या नियमाने}{या नियमांनी}
  तयार होतात. शाखेवरील प्रत्येक पूर्वचिन्हयुक्त
  चिन्हांकित मोडल सूत्रासाठी संबंधित नियम
  एकदाच लावावा लागतो; त्या अनुप्रयोगात
  एकच नवे पूर्वचिन्ह मिळते.
  \iftag{prvBox}{$\TRule{\True}{\Box}$}{}\iftag{notprvBox,notprvDiamond}{}{
    आणि }\iftag{prvDiamond}{$\TRule{\False}{\Diamond}$}{}
  \iftag{notprvBox,notprvDiamond}{हा नियम}{हे नियम}
  अनेकदा लावावे लागू शकतात; पण प्रत्येक
  पूर्वचिन्हयुक्त चिन्हांकित मोडल सूत्रासाठी
  त्याच्या पूर्वचिन्हाचा शाखेवर वापरलेला प्रत्येक
  तात्काळ विस्तार घेऊन एकदाच लावणे पुरते.
  \olref[cpl]{prop:complete-tableau} च्या पुराव्यात
  दाखवल्याप्रमाणे, सांत प्रारंभिक सूत्रसंचाची मोडल
  खोली पूर्वचिन्हांची उंची मर्यादित करते; प्रत्येक
  पूर्वचिन्हाची संततीही सांत असते. म्हणून नवी
  पूर्वचिन्हे सांत संख्येनेच निर्माण होतात. म्हणून
  सांत पायऱ्यांनंतर शाखा बंद होते किंवा
  संपूर्ण होते.

  खुली संपूर्ण शाखा असलेला टॅब्लो
  बांधल्यावर \olref[cpl]{thm:tableau-completeness}
  चा पुरावा मूळ पूर्वचिन्हयुक्त !!{formula}s
  चा संच पूर्ण करणारे स्पष्ट प्रतिमान देतो.
  म्हणून $\Gamma \Entails !A$ असेल, तर
  बंद !!{tableau} अस्तित्वात असतो आणि
  $\Gamma \Proves !A$. याशिवाय, $\Gamma$ सांत
  असेल, तर योग्य पद्धतीने बंद !!{tableau} शोधताना
  त्याऐवजी ‘संपूर्ण’ !!{tableau} मिळाला,
  तर $\Gamma \Entails/ !A$ हे कळते
  आणि प्रतिदृष्टांत प्रतिमान प्रत्यक्ष
  बांधताही येते.
\end{explain}

\iftag{prvBox}{
\begin{ex}
  $\Proves/ \Box(p \lor q) \lif (\Box p \lor \Box
  q)$ हे आपल्याला माहीत आहे. टॅब्लो बांधण्याची
  सुरुवात पुढीलप्रमाणे होते:
  \begin{oltableau}
    [\pFmla{\False}{\Box(p \lor q) \lif (\Box p \lor \Box q)}{1},
      just = \TAss, checked
      [\pFmla{\True}{\Box(p \lor q)}{1},
        just = {\TRule{\False}{\lif}[1]},
        [\pFmla{\False}{\Box p \lor \Box q}{1},
          just = {\TRule{\False}{\lif}[1]}, checked
          [\pFmla{\False}{\Box p}{1},
            just = {\TRule{\False}{\lor}[3]}, checked
            [\pFmla{\False}{\Box q}{1},
              just = {\TRule{\False}{\lor}[3]}, checked
              [\pFmla{\False}{p}{1.1},
                just = {\TRule{\False}{\Box}[4]}, checked
                [\pFmla{\False}{q}{1.2},
                  just = {\TRule{\False}{\Box}[5]}, checked
                ]
              ]
            ]
          ]
        ]
      ]
    ]
  \end{oltableau}
  अर्थात हा !!{tableau} अजून पूर्ण झालेला नाही.
  पुढे खूण नसलेल्या एकमेव ओळीचा विचार करू:
  ओळ~$2$ वरील पूर्वचिन्हयुक्त !!{formula}
  $\sFmla{\True}{\Box(p \lor q)}[1]$.
  पद्धतशीर टॅब्लो बांधताना शाखेवर वापरलेल्या
  $1.n$ या आकाराच्या प्रत्येक पूर्वचिन्हासाठी
  $\TRule{\True}{\Box}$ लावायचा असतो;
  म्हणजे $1.1$ आणि $1.2$ या दोन्हींसाठी:
  \begin{oltableau}
    [\pFmla{\False}{\Box(p \lor q) \lif (\Box p \lor \Box q)}{1},
      just = \TAss, checked
      [\pFmla{\True}{\Box(p \lor q)}{1},
        just = {\TRule{\False}{\lif}[1]},
        [\pFmla{\False}{\Box p \lor \Box q}{1},
          just = {\TRule{\False}{\lif}[1]}, checked
          [\pFmla{\False}{\Box p}{1},
            just = {\TRule{\False}{\lor}[3]}, checked
            [\pFmla{\False}{\Box q}{1},
              just = {\TRule{\False}{\lor}[3]}, checked
              [\pFmla{\False}{p}{1.1},
                just = {\TRule{\False}{\Box}[4]}, checked
                [\pFmla{\False}{q}{1.2},
                  just = {\TRule{\False}{\Box}[5]}, checked
                  [\pFmla{\True}{p \lor q}{1.1},
                    just = {\TRule{\True}{\Box}[2]}
                    [\pFmla{\True}{p \lor q}{1.2},
                      just = {\TRule{\True}{\Box}[2]}
                    ]
                  ]
                ]
              ]
            ]
          ]
        ]
      ]
    ]
  \end{oltableau}
  आता ओळी 2, 8 आणि 9 यांवर खूण नाही.
  पण नवे पूर्वचिन्ह जोडलेले नाही. म्हणून
  ओळी~8 आणि~9 वर
  $\TRule{\True}{\lor}$ लावू; तयार होणाऱ्या
  सर्व शाखांवर, त्या बंद होत नसतील तोवर,
  हाच नियम लागू करू:
  \begin{oltableau}
    [\pFmla{\False}{\Box(p \lor q) \lif (\Box p \lor \Box q)}{1},
      just = \TAss, checked
      [\pFmla{\True}{\Box(p \lor q)}{1},
        just = {\TRule{\False}{\lif}[1]}, checked
        [\pFmla{\False}{\Box p \lor \Box q}{1},
          just = {\TRule{\False}{\lif}[1]}, checked
          [\pFmla{\False}{\Box p}{1},
            just = {\TRule{\False}{\lor}[3]}, checked
            [\pFmla{\False}{\Box q}{1},
              just = {\TRule{\False}{\lor}[3]}, checked
              [\pFmla{\False}{p}{1.1},
                just = {\TRule{\False}{\Box}[4]}, checked
                [\pFmla{\False}{q}{1.2},
                  just = {\TRule{\False}{\Box}[5]}, checked
                  [\pFmla{\True}{p \lor q}{1.1},
                    just = {\TRule{\True}{\Box}[2]}, checked
                    [\pFmla{\True}{p \lor q}{1.2},
                      just = {\TRule{\True}{\Box}[2]}, checked
                      [\pFmla{\True}{p}{1.1},
                        just = {\TRule{\True}{\lor}[8]}, checked, close
                      ]
                      [\pFmla{\True}{q}{1.1},
                        just = {\TRule{\True}{\lor}[8]}, checked
                        [\pFmla{\True}{p}{1.2},
                          just = {\TRule{\True}{\lor}[9]}, checked]
                        [\pFmla{\True}{q}{1.2},
                          just = {\TRule{\True}{\lor}[9]}, checked, close]
                      ]
                    ]
                  ]
                ]
              ]
            ]
          ]
        ]
      ]
    ]
  \end{oltableau}
  आता एकच खुली शाखा उरते; ती संपूर्ण आहे.
  तिच्यावरून प्रतिमान असे ठरवू: जगे
  $W = \{1, 1.1, 1.2\}$ (खुल्या शाखेवर
  येणारी हीच पूर्वचिन्हे), प्राप्यता संबंध
  $R = \{\tuple{1, 1.1}, \tuple{1, 1.2}\}$,
  आणि सत्य-मूल्यांकन $V(p) = \{1.2\}$
  (कारण ओळ~11 वर
  $\sFmla{\True}{p}[1.2]$ आहे) व
  $V(q) = \{1.1\}$ (कारण ओळ~10 वर
  $\sFmla{\True}{q}[1.1]$ आहे).
  \olref{fig:counter-Box} मध्ये हे प्रतिमान
  दाखवले आहे. ते $\Box(p \lor q) \lif
  (\Box p \lor \Box q)$ चे प्रतिदृष्टांत
  प्रतिमान आहे हे तपासता येईल.
  \begin{figure}
  \begin{center}
    \begin{tikzpicture}[modal]
      \node[world] (w1) [label={[align=right]right:\mFalse{p}\\ \mFalse{q}}]{$1$};
      \node[world] (w2) [label={[align=right]right:\mFalse{p}\\ \mTrue{q}},
        above left=of w1]{$1.1$};
      \node[world] (w3) [label={[align=right]right:\mTrue{p}\\ \mFalse{q}},
        above right=of w1] {$1.2$};
      \draw[->] (w1) to (w2);
      \draw[->] (w1) to (w3);
    \end{tikzpicture}
  \end{center}
  \caption{$\Box(p \lor q) \lif (\Box p \lor \Box
    q)$ चे प्रतिदृष्टांत प्रतिमान.}
  \ollabel{fig:counter-Box}
\end{figure}
\end{ex}
}
{
\begin{ex}
  $\Proves/ (\Diamond p \land \Diamond q) \lif \Diamond(p
  \land q)$ हे आपल्याला माहीत आहे. टॅब्लो
  बांधण्याची सुरुवात पुढीलप्रमाणे होते:
  \begin{oltableau}
    [\pFmla{\False}{(\Diamond p \land \Diamond q) \lif \Diamond(p \land q)}{1},
      just = \TAss, checked
      [\pFmla{\True}{\Diamond p \land \Diamond q}{1},
        just = {\TRule{\False}{\lif}[1]}, checked
        [\pFmla{\False}{\Diamond(p \land q)}{1},
          just = {\TRule{\False}{\lif}[1]}
          [\pFmla{\True}{\Diamond p}{1},
            just = {\TRule{\True}{\land}[2]}, checked
            [\pFmla{\True}{\Diamond q}{1},
              just = {\TRule{\True}{\land}[2]}, checked
              [\pFmla{\True}{p}{1.1},
                just = {\TRule{\True}{\Diamond}[4]}, checked
                [\pFmla{\True}{q}{1.2},
                  just = {\TRule{\True}{\Diamond}[5]}, checked
                ]
              ]
            ]
          ]
        ]
      ]
    ]
  \end{oltableau}
  अर्थात हा !!{tableau} अजून पूर्ण झालेला नाही.
  पुढे खूण नसलेल्या एकमेव ओळीचा विचार करू:
  ओळ~$3$ वरील पूर्वचिन्हयुक्त !!{formula}
  $\sFmla{\False}{\Diamond(p \land q)}[1]$.
  पद्धतशीर टॅब्लो बांधताना शाखेवर वापरलेल्या
  $1.n$ या आकाराच्या प्रत्येक पूर्वचिन्हासाठी
  $\TRule{\False}{\Diamond}$ लावायचा असतो;
  म्हणजे $1.1$ आणि $1.2$ या दोन्हींसाठी:
  \begin{oltableau}
    [\pFmla{\False}{(\Diamond p \land \Diamond q) \lif \Diamond(p \land q)}{1},
      just = \TAss, checked
      [\pFmla{\True}{\Diamond p \land \Diamond q}{1},
        just = {\TRule{\False}{\lif}[1]}, checked
        [\pFmla{\False}{\Diamond (p \land q)}{1},
          just = {\TRule{\False}{\lif}[1]}
          [\pFmla{\True}{\Diamond p}{1},
            just = {\TRule{\True}{\land}[2]}, checked
            [\pFmla{\True}{\Diamond q}{1},
              just = {\TRule{\True}{\land}[2]}, checked
              [\pFmla{\True}{p}{1.1},
                just = {\TRule{\True}{\Diamond}[4]}, checked
                [\pFmla{\True}{q}{1.2},
                  just = {\TRule{\True}{\Diamond}[5]}, checked
                  [\pFmla{\False}{p \land q}{1.1},
                    just = {\TRule{\False}{\Diamond}[3]}
                    [\pFmla{\False}{p \land q}{1.2},
                      just = {\TRule{\False}{\Diamond}[3]}
                    ]
                  ]
                ]
              ]
            ]
          ]
        ]
      ]
    ]
  \end{oltableau}
  आता ओळी 3, 8 आणि 9 यांवर खूण नाही.
  पण नवे पूर्वचिन्ह जोडलेले नाही. म्हणून
  ओळी~8 आणि~9 वर
  $\TRule{\False}{\land}$ लावू; तयार होणाऱ्या
  सर्व शाखांवर, त्या बंद होत नसतील तोवर,
  हाच नियम लागू करू:
  \begin{oltableau}
    [\pFmla{\False}{(\Diamond p \land \Diamond q) \lif \Diamond(p \land q)}{1},
      just = \TAss, checked
      [\pFmla{\True}{\Diamond p \land \Diamond q}{1},
        just = {\TRule{\False}{\lif}[1]}, checked
        [\pFmla{\False}{\Diamond(p \land q)}{1},
          just = {\TRule{\False}{\lif}[1]}
          [\pFmla{\True}{\Diamond p}{1},
            just = {\TRule{\True}{\land}[2]}, checked
            [\pFmla{\True}{\Diamond q}{1},
              just = {\TRule{\True}{\land}[2]}, checked
              [\pFmla{\True}{p}{1.1},
                just = {\TRule{\True}{\Diamond}[4]}, checked
                [\pFmla{\True}{q}{1.2},
                  just = {\TRule{\True}{\Diamond}[5]}, checked
                  [\pFmla{\False}{p \land q}{1.1},
                    just = {\TRule{\False}{\Diamond}[3]}, checked
                    [\pFmla{\False}{p \land q}{1.2},
                      just = {\TRule{\False}{\Diamond}[3]}, checked
                      [\pFmla{\False}{p}{1.1},
                        just = {\TRule{\False}{\land}[8]}, checked, close
                      ]
                      [\pFmla{\False}{q}{1.1},
                        just = {\TRule{\False}{\land}[8]}, checked
                        [\pFmla{\False}{p}{1.2},
                          just = {\TRule{\False}{\land}[9]}, checked]
                        [\pFmla{\False}{q}{1.2},
                          just = {\TRule{\False}{\land}[9]}, checked, close]
                      ]
                    ]
                  ]
                ]
              ]
            ]
          ]
        ]
      ]
    ]
  \end{oltableau}
  आता एकच खुली शाखा उरते; ती संपूर्ण आहे.
  तिच्यावरून प्रतिमान असे ठरवू: जगे
  $W = \{1, 1.1, 1.2\}$ (खुल्या शाखेवर
  येणारी हीच पूर्वचिन्हे), प्राप्यता संबंध
  $R = \{\tuple{1, 1.1}, \tuple{1, 1.2}\}$,
  आणि सत्य-मूल्यांकन $V(p) = \{1.1\}$
  (कारण ओळ~6 वर
  $\sFmla{\True}{p}[1.1]$ आहे) व
  $V(q) = \{1.2\}$ (कारण ओळ~7 वर
  $\sFmla{\True}{q}[1.2]$ आहे).
  \olref{fig:counter-Diamond} मध्ये हे प्रतिमान
  दाखवले आहे. ते $(\Diamond p \land \Diamond q)
  \lif \Diamond (p \land q)$ चे प्रतिदृष्टांत
  प्रतिमान आहे हे तपासता येईल.
  \begin{figure}
  \begin{center}
    \begin{tikzpicture}[modal]
      \node[world] (w1) [label={[align=right]right:\mFalse{p}\\ \mFalse{q}}]{$1$};
      \node[world] (w2) [label={[align=right]right:\mTrue{p}\\ \mFalse{q}},
        above left=of w1]{$1.1$};
      \node[world] (w3) [label={[align=right]right:\mFalse{p}\\ \mTrue{q}},
        above right=of w1] {$1.2$};
      \draw[->] (w1) to (w2);
      \draw[->] (w1) to (w3);
    \end{tikzpicture}
  \end{center}
  \caption{$(\Diamond p \land \Diamond q) \lif
    \Diamond (p \land q)$ चे प्रतिदृष्टांत प्रतिमान.}
  \ollabel{fig:counter-Diamond}
\end{figure}
\end{ex}
}

\end{document}
